\documentclass[10pt,aspectratio=169]{beamer}

\usepackage[utf8]{inputenc}
\usepackage[T1]{fontenc}
\usepackage[french]{babel}
\usepackage{amsmath,amssymb,amsfonts}
\usepackage{booktabs}
\usepackage{tikz}
\usepackage{graphicx}
\usepackage{array}
\usepackage{colortbl}
\usepackage{lmodern}
\usepackage{tcolorbox}
\usepackage{pifont}
\usepackage[style=numeric,backend=bibtex]{biblatex}
\usepackage{csquotes}
\usetikzlibrary{positioning,calc,arrows.meta,shapes.geometric,fit}
\tcbuselibrary{skins}
\addbibresource{references.bib}

\definecolor{primary}{RGB}{12,54,106}
\definecolor{secondary}{RGB}{234,241,255}
\definecolor{accent}{RGB}{255,111,0}
\definecolor{darktext}{RGB}{20,33,61}
\definecolor{lighttext}{RGB}{108,117,125}
\definecolor{success}{RGB}{42,157,143}
\definecolor{danger}{RGB}{230,57,70}
\definecolor{warning}{RGB}{244,162,97}

\usefonttheme{professionalfonts}
\setbeamerfont{title}{size=\huge,series=\bfseries}
\setbeamerfont{subtitle}{size=\large,series=\normalfont}
\setbeamerfont{frametitle}{size=\Large,series=\bfseries}
\setbeamercolor{background canvas}{bg=white}
\setbeamercolor{normal text}{fg=darktext}
\setbeamercolor{frametitle}{fg=primary}
\setbeamercolor{title}{fg=primary}
\setbeamercolor{subtitle}{fg=lighttext}
\setbeamercolor{structure}{fg=primary}
\setbeamercolor{item}{fg=accent}

\newtcolorbox{keybox}[1][] {
    colframe=primary!25,
    colback=secondary!35,
    boxrule=0.8pt,
    arc=4pt,
    left=6pt,
    right=6pt,
    top=5pt,
    bottom=5pt,
    #1
}

\newcommand{\cmark}{\textcolor{success}{\ding{51}}}
\newcommand{\xmark}{\textcolor{danger}{\ding{55}}}

\setbeamertemplate{navigation symbols}{}

\begin{document}

\setbeamertemplate{footline}{}
\begin{frame}[plain]
\begin{tikzpicture}[remember picture,overlay]
    \fill[secondary!30] (current page.south west) rectangle (current page.north east);
    \fill[primary!5, opacity=0.85] (current page.north west) -- (current page.south east) -- (current page.south west) -- cycle;
    \fill[accent, opacity=0.08] ($(current page.center)+(3cm,1.2cm)$) circle (4cm);
    \node[anchor=center, align=center] at (current page.center) {
        {\Huge\bfseries\textcolor{primary}{SpeechCoach}}\\[0.3cm]
        {\huge\bfseries\textcolor{darktext}{Réunion d'avancement n\textdegree 3}}\\[0.55cm]
        \begin{tikzpicture}
            \fill[accent] (0,0) rectangle (3.2cm,0.12cm);
        \end{tikzpicture}\\[0.55cm]
        {\Large\itshape\textcolor{lighttext}{Prototype web, moteur analytique et coaching IA local}}\\[1.2cm]
        {\normalsize Hicham Moussaid}\\[0.15cm]
        {\small Master SIE -- ENS Meknès}\\[0.15cm]
        {\small Encadrant : Pr. Abdelaaziz Hessane}
    };
\end{tikzpicture}
\end{frame}

\setbeamertemplate{footline}{
    \begin{tikzpicture}[remember picture,overlay]
        \fill[primary!5] (current page.south west) rectangle ($(current page.south east)+(0,15pt)$);
        \node[anchor=south east, xshift=-14pt, yshift=4pt, text=primary!60, font=\scriptsize\bfseries] at (current page.south east) {
            \insertframenumber~|~\inserttotalframenumber
        };
        \fill[accent] (current page.south west) rectangle ($(current page.south west)+(3pt,15pt)$);
        \node[anchor=south west, xshift=10pt, yshift=4pt, text=primary!60, font=\tiny\bfseries] at (current page.south west) {
            SPEECHCOACH | AVANCEMENT 2026
        };
    \end{tikzpicture}
}

\begin{frame}{Plan}
\begin{columns}[T]
\begin{column}{0.48\textwidth}
\textbf{1. Où en est le projet ?}
\begin{itemize}
    \item objectif produit
    \item état global du prototype
\end{itemize}
\vspace{0.3cm}
\textbf{2. Ce qui a été implémenté}
\begin{itemize}
    \item application web
    \item pipeline analytique
\end{itemize}
\end{column}
\begin{column}{0.48\textwidth}
\textbf{3. Le point difficile : le LLM}
\begin{itemize}
    \item problèmes observés
    \item solutions essayées
\end{itemize}
\vspace{0.3cm}
\textbf{4. Solution retenue et suite}
\begin{itemize}
    \item architecture actuelle
    \item limites et prochaines étapes
\end{itemize}
\end{column}
\end{columns}
\end{frame}

\begin{frame}{Objectif actuel}
\begin{keybox}
\textbf{Idée centrale :} transformer une vidéo de prise de parole en un rapport utile, lisible et actionnable, pas seulement en une liste de métriques.
\end{keybox}
\vspace{0.35cm}
\begin{itemize}
    \item Le projet est passé d'un pipeline technique à un \textbf{prototype web complet}.
    \item Le cœur reste \textbf{multimodal} : audio, regard, gestuelle, qualité visuelle.
    \item Le but de cette phase était surtout de rendre le système \textbf{présentable, stable et exploitable} pour un vrai usage.
\end{itemize}
\vspace{0.3cm}
% \begin{tcolorbox}[colback=success!6,colframe=success!35]
% \small \textbf{Message cle a dire :} cette phase n'a pas seulement ajoute des fonctions, elle a transforme le prototype en produit utilisable.
% \end{tcolorbox}
\end{frame}

\begin{frame}{Architecture actuelle du prototype}
\centering
\begin{tikzpicture}[
    node distance=1.3cm and 1.5cm,
    box/.style={draw, rounded corners, align=center, minimum width=3.2cm, minimum height=1.05cm, fill=blue!5, font=\scriptsize},
    service/.style={draw, rounded corners, align=center, minimum width=3.2cm, minimum height=1.05cm, fill=green!5, font=\scriptsize},
    storage/.style={draw, rounded corners, align=center, minimum width=3.2cm, minimum height=1.05cm, fill=orange!8, font=\scriptsize},
    arrow/.style={-{Latex[length=2.3mm]}, thick}
]
\node[box] (frontend) {Frontend web\\Next.js / React / TypeScript};
\node[box, right=of frontend] (api) {Backend API\\FastAPI / SQLModel};
\node[service, right=of api] (ollama) {Coaching IA local\\Ollama + GGUF};
\node[storage, below=of frontend] (db) {Base de donnees\\utilisateurs / profils / sessions};
\node[service, below=of api] (celery) {Traitement asynchrone\\Celery + Redis};
\node[service, below=of ollama] (engine) {Moteur analytique\\FFmpeg, Whisper, MediaPipe, OpenCV};
\node[storage, below=of celery] (storage) {Stockage local\\uploads / rapports};
\draw[arrow] (frontend) -- (api);
\draw[arrow] (api) -- (celery);
\draw[arrow] (celery) -- (engine);
\draw[arrow] (engine) -- (ollama);
\draw[arrow] (api) -- (db);
\draw[arrow] (celery) -- (storage);
\draw[arrow] (api) -- (ollama);
\end{tikzpicture}
\vspace{0.07cm}
\begin{itemize}
    \item Le calcul analytique et le coaching sont \textbf{separes} pour garder une execution stable.
    \item Le LLM ne decide pas seul : il \textbf{reformule} une base deja produite par le moteur deterministe.
\end{itemize}
\end{frame}

\begin{frame}{Ce qui a ete implemente}
\begin{columns}[T]
\begin{column}{0.49\textwidth}
\begin{tcolorbox}[title=Application web, colframe=primary, colback=primary!3]
\small
\begin{itemize}
    \item inscription et connexion
    \item gestion du profil et des parametres
    \item dashboard et historique des analyses
    \item export PDF et Markdown
\end{itemize}
\end{tcolorbox}
\begin{tcolorbox}[title=Studio, colframe=accent, colback=accent!3]
\small
\begin{itemize}
    \item import d'une video
    \item enregistrement direct webcam/micro
    \item suivi du traitement et consultation du rapport
\end{itemize}
\end{tcolorbox}
\end{column}
\begin{column}{0.49\textwidth}
\begin{tcolorbox}[title=Moteur analytique, colframe=success, colback=success!3]
\small
\begin{itemize}
    \item transcription et debit de parole
    \item pauses, fillers, energie
    \item regard camera, mains visibles, intensite gestuelle
    \item qualite visuelle et cadrage
\end{itemize}
\end{tcolorbox}
\begin{tcolorbox}[title=Restitution, colframe=warning!70!black, colback=warning!8]
\small
\begin{itemize}
    \item score global et scores par axe
    \item points forts / axes de progression
    \item plan d'exercice et transcription
\end{itemize}
\end{tcolorbox}
\end{column}
\end{columns}
\end{frame}

\begin{frame}{Pourquoi le LLM a été difficile}
\begin{keybox}[colframe=danger!35,colback=danger!4]
\textbf{Le vrai problème :} demander à un LLM local de produire \textbf{trois textes courts, justes, sobres et en français propre} est plus difficile que de générer un long paragraphe libre.
\end{keybox}
\vspace{0.35cm}
\begin{itemize}
    \item Erreurs observées : hallucinations, ton générique, fuite de prompt, JSON invalide, confusion entre force et faiblesse.
    \item Le contexte \textbf{RAG} aidait à enrichir l'appui pédagogique, mais pouvait aussi perturber le \textbf{bilan court}.
    \item Le problème n'était donc pas seulement le choix du modèle, mais aussi \textbf{l'architecture de génération}. 
\end{itemize}
\vspace{0.3cm}
\begin{tcolorbox}[colback=secondary!45,colframe=primary!25]
\small \textbf{Le LLM était utile, mais pas assez fiable pour devenir la source principale de vérité.}
\end{tcolorbox}
\end{frame}

\begin{frame}{Solutions essayées pendant cette phase}
\begin{columns}[T]
\begin{column}{0.49\textwidth}
\begin{itemize}
    \item comparaison de plusieurs modèles locaux : Llama, Qwen, Mistral
    \item prompt engineering plus strict
    \item validation JSON et fallbacks déterministes
    \item ajout du RAG pour donner plus de contexte pédagogique au coaching
\end{itemize}
\end{column}
\begin{column}{0.49\textwidth}
\begin{itemize}
    \item fine-tuning supervisé avec \textbf{QLoRA}
    \item dataset construit en \textbf{input/target}, puis converti en \textbf{messages}
    \item essais sur plusieurs tailles de modèles
    \item tests de déploiement local en CPU
\end{itemize}
\end{column}
\end{columns}
\vspace{0.35cm}
\begin{tcolorbox}[colback=accent!8,colframe=accent!45]
\small \textbf{Ce qu'on a appris :} le fine-tuning a amélioré le format et le style, mais le problème du déploiement local restait entier avec l'approche base model + adapter.
\end{tcolorbox}
\end{frame}

\begin{frame}{Dataset de fine-tuning : construction et nettoyage}
\vspace{-0.35cm}

\begin{columns}[T]
\begin{column}{0.49\textwidth}
\begin{tcolorbox}[title=Comment le dataset a été construit, colframe=primary, colback=primary!3]
\small
\begin{itemize}
    \item point de départ : cas réels issus des rapports du projet
    \item ajout de cas synthétiques pour couvrir plus de scénarios
    \item objectif : apprendre le ton, la structure et le format des 3 champs courts
\end{itemize}
\end{tcolorbox}
\begin{tcolorbox}[title=Nettoyage réalisé, colframe=warning!60!black, colback=warning!8]
\small
\begin{itemize}
    \item suppression des doublons
    \item suppression des cas incohérents ou trop artificiels
    \item recentrage sur les catégories réellement utilisées par le moteur
    \item version finale compacte et exploitable
\end{itemize}
\end{tcolorbox}
\end{column}
\begin{column}{0.49\textwidth}
\begin{tcolorbox}[title=Structure des données, colframe=success, colback=success!4]
\small
\textbf{Étape 1 : format lisible pour nous}
\begin{itemize}
    \item \texttt{input} : scores, forces, faiblesses, top recommendation
    \item \texttt{target} : bilan, priorité, encouragement
\end{itemize}
\vspace{0.15cm}
\textbf{Étape 2 : conversion pour l'entraînement}
\begin{itemize}
    \item ajout d'un message system
    \item transformation en format \texttt{messages}
    \item compatibilite avec QLoRA / Transformers
\end{itemize}
\end{tcolorbox}
\end{column}
\end{columns}
\vspace{0.25cm}
\begin{tcolorbox}[colback=secondary!45,colframe=primary!25]
\small \textbf{Message clé :} le dataset n'a pas été un export brut ; il a été construit, nettoyé puis converti pour rester cohérent avec le vrai comportement de SpeechCoach.
\end{tcolorbox}
\end{frame}

\begin{frame}{Architecture retenue pour le coaching IA}
\begin{columns}[T]
\begin{column}{0.48\textwidth}
\textbf{Avant}
\begin{itemize}
    \item modele de base + adapter LoRA separes
    \item chargement lourd dans le worker Python
    \item temps longs et risque de crash memoire
\end{itemize}
\vspace{0.2cm}
\textbf{Maintenant}
\begin{itemize}
    \item fusion des poids
    \item quantification GGUF
    \item service local via Ollama
\end{itemize}
\end{column}
\begin{column}{0.48\textwidth}
\begin{tcolorbox}[title=Architecture actuelle, colframe=success, colback=success!4]
\small
\begin{itemize}
    \item scoring et priorites : \textbf{deterministes}
    \item exercices : \textbf{base de connaissances + mapping deterministe}
    \item bilan / priorite / encouragement : \textbf{LLM local enrichi par RAG}
    \item nettoyage et fallback : \textbf{deterministes}
\end{itemize}
\end{tcolorbox}
\end{column}
\end{columns}
\vspace{0.3cm}
\begin{tcolorbox}[colback=primary!4,colframe=primary!25]
\small \textbf on n'a pas remplace le systeme par un LLM ; on a garde un coeur fiable et on a utilise le LLM seulement la ou il apporte une vraie valeur.
\end{tcolorbox}
\end{frame}

\begin{frame}{État actuel du prototype}
\begin{columns}[T]
\begin{column}{0.49\textwidth}
\begin{tcolorbox}[title=Points solides, colframe=success, colback=success!4]
\small
\begin{itemize}
    \item application web deja exploitable
    \item pipeline analytique stable
    \item rapport plus lisible pour l'utilisateur
    \item coaching IA local fonctionnel
\end{itemize}
\end{tcolorbox}
\end{column}
\begin{column}{0.49\textwidth}
\begin{tcolorbox}[title=Points à finaliser, colframe=warning!60!black, colback=warning!8]
\small
\begin{itemize}
    \item polissage final du texte généré
    \item adaptation aux différents appareils
    \item plan de progression encore trop statique
    \item interface et responsive à consolider
\end{itemize}
\end{tcolorbox}
\end{column}
\end{columns}
\vspace{0.35cm}
\begin{keybox}
\textbf{Le projet est déjà démontrable de bout en bout ; la suite porte surtout sur le polissage, l'évaluation et la consolidation.}
\end{keybox}
\end{frame}

\begin{frame}{Prochaines étapes}
\begin{itemize}
    \item adapter le système à différents appareils : smartphone, tablette, PC, avec recalibrage des signaux comme pitch, yaw et regard caméra
    \item faire évoluer le plan de progression, aujourd'hui encore trop statique, vers un format plus court et plus personnalisé que le schéma fixe sur 7 jours
    \item rendre l'interface vraiment responsive et ajouter les briques produit restantes : vérification e-mail, derniers raffinements des paramètres et du parcours utilisateur
    \item enrichir la restitution avec d'autres graphiques et un meilleur niveau de finition visuelle
\end{itemize}
\vspace{0.4cm}
\begin{tcolorbox}[colback=secondary!45,colframe=primary!25]
\small \textbf{La prochaine phase ne consiste pas à changer l'idée du projet, mais à rendre le système plus adaptable, plus personnalisable et plus propre côté produit.}
\end{tcolorbox}
\end{frame}

\setbeamertemplate{footline}{}
\begin{frame}[plain]
\begin{tikzpicture}[remember picture, overlay]
    \fill[primary] (current page.south west) rectangle (current page.north east);
    \node[text=white, align=center] at (current page.center) {
        {\Huge\bfseries Merci de votre attention}\\[0.45cm]
        {\Large Questions et demonstration}\\[0.9cm]
        \begin{tikzpicture}
            \fill[accent] (0,0) rectangle (4cm,0.14cm);
        \end{tikzpicture}
    };
\end{tikzpicture}
\end{frame}
% --- SLIDE 14: RÉFÉRENCES ---
\begin{frame}[allowframebreaks]{Références Bibliographiques}
    \nocite{*}
    \printbibliography[heading=none]
\end{frame}

\end{document}
